% =============================================================================
% preamble.tex — Category Theory Book
% Input this file from main.tex after \documentclass and before \begin{document}
% =============================================================================

% --- Packages: Mathematics ---
\usepackage{amsmath}
\usepackage{amssymb}
\usepackage{amsthm}
\usepackage{mathtools}

% --- Packages: Diagrams ---
\usepackage{tikz-cd}

% --- Packages: Color and Boxes ---
\usepackage{xcolor}
\usepackage{tcolorbox}
\tcbuselibrary{skins}

% --- Packages: Page Layout ---
\usepackage[margin=1.25in]{geometry}

% --- Packages: References and Links ---
\usepackage[
  colorlinks=true,
  linkcolor=blue!60!black,
  citecolor=green!50!black,
  urlcolor=blue
]{hyperref}
\usepackage{cleveref}  % must come AFTER hyperref

% --- Packages: Lists ---
\usepackage{enumitem}

% --- Packages: Bibliography ---
\usepackage[backend=biber, style=alphabetic]{biblatex}

% --- Packages: Typography and Tables ---
\usepackage{microtype}
\usepackage{booktabs}
\usepackage{graphicx}

% =============================================================================
% --- TikZ Libraries ---
% =============================================================================

\usetikzlibrary{
  arrows.meta,
  decorations.pathmorphing,
  positioning,
  calc
}

% --- Commutative Diagram Defaults ---
\tikzset{
  commutative diagrams/.cd,
  arrow style=tikz,
  diagrams={>=Latex}
}

% =============================================================================
% --- Theorem Environments ---
% =============================================================================

\theoremstyle{plain}
\newtheorem{theorem}{Theorem}[chapter]
\newtheorem{proposition}[theorem]{Proposition}
\newtheorem{lemma}[theorem]{Lemma}
\newtheorem{corollary}[theorem]{Corollary}

\theoremstyle{definition}
\newtheorem{definition}[theorem]{Definition}
\newtheorem{example}[theorem]{Example}

\theoremstyle{remark}
\newtheorem{remark}[theorem]{Remark}
\newtheorem{warning}[theorem]{Warning}
\newtheorem{exercise}[theorem]{Exercise}

% =============================================================================
% --- Math Macros ---
% =============================================================================

% Generic category name (bold)
\newcommand{\cat}[1]{\mathbf{#1}}

% Standard categories
\newcommand{\Set}{\mathbf{Set}}
\newcommand{\Grp}{\mathbf{Grp}}
\newcommand{\Top}{\mathbf{Top}}
\newcommand{\Vect}{\mathbf{Vect}}
\newcommand{\Ring}{\mathbf{Ring}}
\newcommand{\Pos}{\mathbf{Pos}}
\newcommand{\Cat}{\mathbf{Cat}}
\newcommand{\Ab}{\mathbf{Ab}}
\newcommand{\Man}{\mathbf{Man}}

% Opposite category
\newcommand{\op}{^{\mathrm{op}}}

% Hom, End, Aut
\newcommand{\Hom}{\operatorname{Hom}}
\newcommand{\End}{\operatorname{End}}
\newcommand{\Aut}{\operatorname{Aut}}

% Identity morphisms
\newcommand{\id}{\mathrm{id}}
\newcommand{\Id}{\mathrm{Id}}

% Yoneda embedding
\newcommand{\yo}{\mathbf{y}}

% Adjunction symbol
\newcommand{\adj}{\mathrel{\dashv}}

% Kan extensions
\newcommand{\Lan}{\operatorname{Lan}}
\newcommand{\Ran}{\operatorname{Ran}}

% Colimit
\newcommand{\colim}{\operatorname{colim}}

% Natural transformations and functor categories
\newcommand{\Nat}{\operatorname{Nat}}
\newcommand{\Fun}{\operatorname{Fun}}

% Objects and morphisms of a category
\newcommand{\ob}{\operatorname{ob}}
\newcommand{\mor}{\operatorname{mor}}

% Domain and codomain
\newcommand{\cod}{\operatorname{cod}}
\newcommand{\dom}{\operatorname{dom}}

% Isomorphism arrow
\newcommand{\iso}{\xrightarrow{\sim}}

% Number systems
\newcommand{\NN}{\mathbb{N}}
\newcommand{\ZZ}{\mathbb{Z}}
\newcommand{\QQ}{\mathbb{Q}}
\newcommand{\RR}{\mathbb{R}}
\newcommand{\CC}{\mathbb{C}}

% Tensor product shorthand
\newcommand{\tensor}{\otimes}

% Blank slot / placeholder in a functor
\newcommand{\blank}{{-}}

% =============================================================================
% --- tcolorbox Styles ---
% =============================================================================

\tcbset{
  conceptbox/.style={
    enhanced,
    colback=blue!5!white,
    colframe=blue!40!black,
    fonttitle=\bfseries,
    boxrule=0.6pt,
    arc=4pt,
    left=6pt,
    right=6pt,
    top=4pt,
    bottom=4pt
  }
}
